\scanpage{8}
\deleted{\uline{Théorème} 5.1. bis \quad Soient $G$ un $S$-préschéma en groupes de présentation finie, $H_1$ et $H_2$ deux sous-préschémas en groupes de type multiplicatif, \deleted{\unsure{sur $S$ tel que $H$|s $S$ tel que $H$}} Soit $s\quad S$ \deleted{\unsure{et su}} \deleted{\unsure{y}}tel que $H_{1\,s}=H_{2\,s}$, et supposons $H_1$ central. Alors il existe un voisinage ouvert $U$ de $s$ tel que $H_{1\,U}=H_{2\,U}$.}

\deleted{\uline{Corol}laire 5.2. bis}

\uline{Corollaire} 5.3. \quad Soient $u,H,G$, comme dans 5.2. Supposons de plus $G$ \uline{séparé} sur $S$. Soit $U$ l'ensemble des $s\quad S$ tels que $u_s$ soit l'homomorphisme unité. Alors $U$ est une partie à la fois ouverte et fermée de $S$, et $u_U:H_U\quad G_U$ est l'homomorphisme unité.

Il reste à prouver seulement que $U$ est fermé. Or \deleted{pour tout entier} $n\quad 0$, soit $ {}_nu:{}_nH\quad G$ l'homomorphisme induit, et $U_n$ l'ensemble des $s\quad S$ tels que $ {}_nu_s$ soit l'homomorphisme unité. En vertu de 4.6. \unsure{appliqué|applique} aux fibres $H_s$, $U$ est l'intersection des $U_n$, donc on est ramené à prou\deleted{ver que ces derniers sont fermés, donc on peut supposer $G$ \uline{fini} sur $S$.} soit $K=\operatorname{Ker}u$, comme $H$ est séparé sur $S$, c'est un sous-préschéma \uline{fermé} de $H$, et $U$ est l'ensemble des $s\quad S$ tels que $K_s=$\deleted{\unsure{encore}}$H_s$ . On voit alors \deleted{\unsure{facistem|facislem|facilemen}} faci\deleted{\unsure{m}}lement que $U$ est fermé, par exemple comme application de Exp V,6.1. ,ou en \deleted{\unsure{se ramenant}} notant qu'on peut supposer $S$ réduit et $U$ dense dans $S$, donc schématiquement dense dans $S$, ce qui implique $H_U$ schématiquement dense dans $H$ puisque $H$ est plat sur $S$, et comme $u$ et l'homomorphisme unité $e$ coïncident sur $H_U$, ils coincident sur $H$.

La démonstration est la même que pour 5.1. , en rempla,cant 3.4. par 3.4. bis et 4.7. a) par 4.7. b).

\deleted{\uline{Corollaire} 5.3. bis}
% archive sheet number: 7

